\subsection{One accepted main path with protected challenge boundaries}\label{sec:global:checkpoints}
The only new proof mechanism in this chapter is an
\emph{extract-and-restore} operation at the end of an
already accepted sequential OR proof. It does not change
the verifier messages on the main execution path.
All extra executions run from saved classical prover
states, with the \emph{same} prover random tape, and their
messages are discarded before continuing the main path.
This is a stand-alone, classical interactive reduction;
the external KBFold oracle is explicitly separated.

\begin{definition}[Three outer epochs]\label{def:global:epochs}
A \emph{two-challenge execution} is divided into
three ordered epochs:
\begin{enumerate}[label=\textup{(\roman*)}]
\item \emph{Epoch 0}, before the first fresh independent
outer challenge $\theta\leftarrow L^\ell$:
the hidden witness/packed-column bit coordinates
and the hidden field variables $p,p_{\rm in},c$
are registered with mandatory bitness proofs;
the public masked values $m,y'$ are published
and fixed. Every variable entering the first
discrepancy is thus fixed before $\theta$.
\item \emph{Epoch 1}, between $\theta$ and the second
fresh independent outer challenge $z\leftarrow L$:
all additional hidden variables $t,\rho,\sigma$
and the public trace vector $U$ are fixed, with
their mandatory bitness proofs completed.
\item \emph{Epoch 2}, after $z$: all remaining
prescribed public-coefficient affine parity
proofs and authenticated-oracle checks are completed.
An affine parity proof may also occur in an earlier
epoch once its public coefficients are available.
\end{enumerate}
No commitment or field variable from a later
epoch may replace an already registered value
used in an earlier discrepancy polynomial.
\end{definition}

\begin{definition}[Sequential OR blocks]\label{def:global:block}
An \emph{OR block} consists of one public fixed
statement $(C,S)$ and $\kappa\ge1$ sequential
one-bit OR mini-proofs of
Definition~\ref{def:mv:repeat}, with one
fresh uniform verifier bit after each first
message. The prover's behavior, including
later first messages in this block, may depend
on earlier bits in that same block.
The verifier requires all $\kappa$ mini-proofs
to accept before continuing.
There are at most $J$ OR blocks on any
accepting main path; $J$ is a public
polynomial bound.
\end{definition}

\begin{definition}[Checkpoint-replay access]\label{def:global:access}
Immediately before each OR block, a classical
black-box reduction may save the prover's
complete machine state and random tape.
It can execute alternative sequences of
inner verifier bits from that state and
then \emph{restore the exact state at the
end of the already sampled main block}.
The reduction does not send alternative
transcripts to the main verifier.
It cannot rewind a quantum prover,
cannot change an outer challenge already
sampled, and cannot reset an external
stateful oracle unless that capability is
included in the explicit external model.
\end{definition}

\begin{definition}[Main and side executions]\label{def:global:main}
At every OR block the reduction first runs
one \emph{main} execution using the usual
independent uniform verifier bits. If the
main block rejects, the entire main
execution rejects. If it accepts, the
reduction may run \emph{side} executions
from the saved pre-block checkpoint,
using new independent uniform inner bits,
before restoring the accepted main
post-block state and proceeding.
Every later main inner or outer challenge
is sampled exactly as in the unmodified
honest-verifier protocol.
\end{definition}

\begin{lemma}[Replay preserves the main-path law]\label{lem:global:main-law}
For every causal randomized classical prover
satisfying Definition~\ref{def:global:access},
the public main transcript, including both outer
challenges and all inner verifier bits, has
exactly the same distribution with or without
any number of checkpoint-replay side executions.
\end{lemma}
\begin{proof}
Fix the prover's random tape and induct over
the ordered main protocol stages.
At a stage with no OR block, both executions
send the same messages and sample the same
distribution of verifier challenges.
At an OR block, the main challenge sequence
and the resulting accepting or rejecting
post-block state are sampled first, just
as in the unmodified protocol.
Every side execution begins from the
saved pre-block state, and the reduction
restores the fixed sampled main post-block
state before the next main stage.
Thus the side runs do not alter any
main message, state or verifier coin.
Induction gives equality of conditional
main-transcript laws for the fixed
prover tape. Averaging over that tape
proves the lemma.
\end{proof}

\begin{lemma}[Protected outer challenges stay independent]\label{lem:global:outer-uniform}
Suppose $\theta$ is sampled by the honest
verifier only after all Epoch 0 extraction
attempts, and $z$ only after all Epoch 1
extraction attempts. Conditional on the
entire state of the reduction immediately
before either sampling, that outer
challenge has its prescribed fresh
uniform law.
\end{lemma}
\begin{proof}
The verifier samples its outer challenge
from fresh independent randomness, which
has not been used in either the main
preceding transcript or any side replay.
The reduction's pre-challenge state is a
function of earlier prover states, inner
bits and replay coins only. Independence
of the fresh random sample from that
sigma-field yields the uniform
conditional law.
\end{proof}

\begin{remark}[Essential scheduling rule]
The bitness proof for a value that enters
the first challenge-dependent discrepancy
must finish in Epoch 0; a bitness proof for
an auxiliary residual value must finish in
Epoch 1. Simply committing to a Pedersen
group element before a challenge and
deferring its \emph{knowledge proof}
until afterward does not meet this
extractor schedule.
\end{remark}
